\documentclass[11pt]{article}
\usepackage[letterpaper,margin=0.72in,top=0.82in,bottom=0.72in,headheight=22pt]{geometry}
\usepackage[T1]{fontenc}
\usepackage{lmodern}
\usepackage{microtype}
\usepackage[table]{xcolor}
\usepackage{amsmath,amssymb,mathtools}
\usepackage{array,longtable,booktabs,tabularx}
\usepackage[inline]{enumitem}
\usepackage[most]{tcolorbox}
\usepackage{fancyhdr}
\usepackage{hyperref}

\definecolor{Primary}{HTML}{17324D}
\definecolor{Accent}{HTML}{2C7A7B}
\definecolor{CueTint}{HTML}{EAF1F5}
\definecolor{NoteTint}{HTML}{F8FAFB}
\definecolor{Rule}{HTML}{B8C6CF}
\hypersetup{hidelinks,pdftitle={Chapter 3: Exact Matching: A Deeper Look at Classical Methods},pdfsubject={Cornell notes}}
\setlist[itemize]{leftmargin=1.25em,itemsep=1pt,topsep=1.5pt,parsep=0pt}
\setlength{\parindent}{0pt}
\setlength{\parskip}{3pt}
\renewcommand{\arraystretch}{1.16}
\emergencystretch=2em

\pagestyle{fancy}
\fancyhf{}
\fancyhead[L]{\sffamily\small\color{Primary} Chapter 3}
\fancyhead[R]{\sffamily\small\color{Primary} Advanced exact matching}
\fancyfoot[C]{\sffamily\small\color{Primary}\thepage}
\renewcommand{\headrulewidth}{0.5pt}
\renewcommand{\headrule}{\hbox to\headwidth{\color{Accent}\leaders\hrule height \headrulewidth\hfill}}

\newcolumntype{C}{>{\raggedright\arraybackslash}p{0.275\textwidth}}
\newcolumntype{N}{>{\raggedright\arraybackslash}p{0.655\textwidth}}
\newcommand{\cnrow}[2]{%
  \cellcolor{CueTint}\textbf{\sffamily\color{Primary}#1} &
  \cellcolor{NoteTint}#2 \\[4pt]\arrayrulecolor{Rule}\hline}

\begin{document}

\begin{tcolorbox}[
  enhanced,breakable,colback=Primary!4,colframe=Primary,boxrule=0.9pt,
  arc=2mm,left=3mm,right=3mm,top=2.5mm,bottom=2.5mm]
{\sffamily\color{Accent}\bfseries CHAPTER 3}\\[-1pt]
{\sffamily\LARGE\bfseries\color{Primary} Exact Matching: A Deeper Look at Classical Methods}

\vspace{2pt}
\textbf{Purpose.} Refine the analysis of classical matchers and extend single-pattern failure transitions to dictionaries, wildcard patterns, grids, and regular languages.

\textbf{Learning objectives}
\begin{itemize}
  \item Explain how remembered match lengths prevent redundant Boyer-Moore comparisons.
  \item Understand the shape of a linear comparison-bound proof.
  \item Build and use failure links in a keyword tree.
  \item Recognize reductions from wildcard and two-dimensional matching to exact set matching.
\end{itemize}

\textbf{Prerequisites.} Boyer-Moore, KMP, tries, finite automata, and amortized charging arguments.
\end{tcolorbox}

\begin{longtable}{@{}C!{\color{Accent}\vrule width 0.8pt}N@{}}
\rowcolor{Primary}
\color{white}\sffamily\bfseries Cue / question &
\color{white}\sffamily\bfseries Notes \\ \hline
\endfirsthead
\rowcolor{Primary}
\color{white}\sffamily\bfseries Cue / question &
\color{white}\sffamily\bfseries Notes (continued) \\ \hline
\endhead
\cnrow{\S3.1: How can Boyer-Moore remember work?}{A refined variant records how many characters were previously verified at text positions near the current alignment. During a later right-to-left scan, stored lengths can skip a block or certify a mismatch boundary. The proof partitions work into phases and charges new comparisons to progress in the text.}
\cnrow{What is the mini-profile of the skip variant?}{\textbf{Problem:} exact matching with Boyer-Moore-style skips. \textbf{State:} shift tables plus remembered matched-suffix lengths. \textbf{Invariant:} skipped characters were already certified in a compatible alignment. \textbf{Bound:} linear worst-case search while retaining large practical shifts. \textbf{Risk:} stale memory must never be reused outside the alignment conditions that justified it.}
\cnrow{\S3.2: What does Cole's analysis establish?}{The original Boyer-Moore rules admit a linear worst-case comparison bound, including cases with occurrences. The analysis groups comparisons around shifts and uses the structure of matched suffixes and periods to show that costly alignments force enough forward progress. The significance is a true constant-times-$n$ bound, not merely good average behavior.}
\cnrow{\S3.3: What was the original KMP preprocessing idea?}{Compute failure shifts directly from borders of successive prefixes. The easy case extends the previous border when the next characters agree; otherwise follow shorter failure links until extension becomes possible. Optimized shifts can bypass a state that would immediately compare the same mismatching character again.}
\cnrow{\S3.4: How does a keyword tree represent many patterns?}{Insert every pattern into a trie. Trie edges consume characters and terminal nodes identify outputs. A text scan follows trie edges as far as possible, but naive restart after a failure would repeat work.}
\cnrow{What are Aho-Corasick failure links?}{For a node spelling string $x$, its failure link points to the deepest node spelling a proper suffix of $x$ that is also a trie prefix. Build links breadth first so a parent's link is already known. Output information propagates along failure links, allowing suffix patterns to be reported.}
\cnrow{Why is dictionary matching linear?}{Each text character causes a transition or a sequence of failure steps followed by a transition. Depth increases are tied to consumed characters and failure steps decrease depth, yielding linear total transition work. With complete transition tables, each character can be processed directly. Runtime also includes the number of reported matches.}
\cnrow{\S3.5: Where does exact set matching apply?}{A library of known biological strings is a direct dictionary. A wildcard pattern can be split into maximal literal pieces; exact occurrences of those pieces generate candidate alignments that are then combined. A two-dimensional pattern can encode rows and match the resulting row-name sequence, provided row occurrences and alignment constraints are tracked.}
\cnrow{\S3.6: How do regular expressions broaden the model?}{A regular expression denotes a language rather than a fixed word. Compile it to a finite automaton and simulate reachable states while reading the text. Deterministic simulation gives one active state; nondeterministic simulation maintains a set. The accepted language and desired substring semantics determine initialization and reporting.}
\cnrow{What should be compared across methods?}{Record preprocessing time, automaton or trie size, text-scan time, output size, alphabet dependence, and whether overlapping results are reported. A method linear in text length may still be dominated by a large dictionary, transition table, or output set.}
\cnrow{\S3.7: What should practice emphasize?}{Trace failure-link construction, prove that a failure target is the longest usable suffix, and design reductions that preserve both occurrence positions and total complexity. Test nested patterns such as \texttt{he}, \texttt{she}, and \texttt{hers} to expose output propagation.}
\end{longtable}


\begin{tcolorbox}[enhanced,breakable,title={Chapter synthesis},
  colback=Accent!5,colframe=Accent,fonttitle=\small\sffamily\bfseries,fontupper=\footnotesize,
  boxsep=0.6mm,left=1.5mm,right=1.5mm,top=1mm,bottom=1mm,before skip=3pt,after skip=3pt]
The deeper view replaces folklore about practical speed with invariants and charging proofs. Remembered matches make Boyer-Moore comparisons reusable, direct border computation explains KMP's original form, and failure links lift the same suffix-prefix logic from a path to a trie. This creates a common foundation for dictionaries and richer pattern models.
\end{tcolorbox}

\begin{tcolorbox}[enhanced,breakable,title={Key takeaways},
  colback=white,colframe=Primary,fonttitle=\small\sffamily\bfseries,fontupper=\footnotesize,
  boxsep=0.6mm,left=1.5mm,right=1.5mm,top=1mm,bottom=1mm,before skip=3pt,after skip=3pt]
\begin{itemize}
  \item Skipping requires a certificate that earlier comparisons remain applicable.
  \item Worst-case comparison proofs must link expensive work to irreversible progress.
  \item Trie failure links are the many-pattern analogue of KMP fallback links.
  \item Dictionary search time is output sensitive.
  \item Wildcard and grid matching need reductions that preserve alignment coordinates.
\end{itemize}
\end{tcolorbox}

\begin{tcolorbox}[enhanced,breakable,title={Algorithm selection / when to use this},
  colback=Primary!3,colframe=Primary,fonttitle=\small\sffamily\bfseries,fontupper=\footnotesize,
  boxsep=0.6mm,left=1.5mm,right=1.5mm,top=1mm,bottom=1mm,before skip=3pt,after skip=3pt]
\begin{itemize}
  \item Use remembered-match Boyer-Moore variants when both skipping and a simple linear bound matter.
  \item Use Aho-Corasick for simultaneous search of many fixed strings.
  \item Use finite automata when the pattern denotes a regular language rather than a finite dictionary.
\end{itemize}
\end{tcolorbox}

\begin{tcolorbox}[enhanced,breakable,title={Self-test questions},
  colback=white,colframe=Accent,fonttitle=\small\sffamily\bfseries,fontupper=\footnotesize,
  boxsep=0.6mm,left=1.5mm,right=1.5mm,top=1mm,bottom=1mm,before skip=3pt,after skip=3pt]
\begin{itemize}
  \item What condition makes a remembered match length safe to reuse?
  \item Define the failure link of a keyword-tree node.
  \item Why is breadth-first order convenient for constructing failure links?
\end{itemize}
\end{tcolorbox}

{\footnotesize
\textbf{Terms and notation to remember.} skip array, border optimization, keyword tree, failure link, output link, finite automaton.

\textbf{Connections.} Arithmetic filters and suffix indexes provide the next two alternatives to comparison-based matching.
}

\end{document}
